\documentclass[12 pt, letterpaper]{hmcpset}
\usepackage{hw-preamble}

\name{Jayden Thadani}
\class{Analysis}
\assignment{Chapter 2 exercises --- Basic Topology}
\duedate{}

\begin{document}

\begin{problem}[1.]
	Prove that the empty set is a subset of every set.
\end{problem}

\pagebreak

\begin{problem}[2.]
	A complex number $z$ is said to be \textit{algebraic} if there are integers $a_{0}, \dots, a_{n}$ not all zero, such that
	\[
		a_{0} z^{n} + a_{1} z^{n - 1} + \dots + a_{n - 1} z + a_{n} = 0.
	\]
	Prove that the set of all algebraic numbers is countable.
\end{problem}

\pagebreak

\begin{problem}[3.]
	Prove that there exist real numbers which are not algebraic.
\end{problem}

\pagebreak

\begin{problem}[4.]
	Is the set of all irrational real numbers countable?
\end{problem}

\pagebreak

\begin{problem}[5.]
	Construct a bounded set of real numbers with exactly three limit points.
\end{problem}

\pagebreak

\begin{problem}[6.]
	Let $E'$ be the set of all limit points of a set $E$. Prove that $E'$ is closed. Prove that $E$ and $\overline{E}$ have the same limit points. Do $E$ and $E'$ always have the same limit points?
\end{problem}

\pagebreak

\begin{problem}[7.]
	Let $A_{1}, A_{2}, A_{3}, \dots$ be subsets of a metric space.
	\begin{enumerate}
		\item If $B_{n} = \bigcup_{i = 1}^{n} A_{i}$, prove that $\overline{B_{n}} = \bigcup_{i = 1}^{n} \overline{A_{i}}$, for $n \in \mathbb{N}$.
		\item If $B = \bigcup_{i \in \mathbb{N}} A_{i}$, prove that $\overline{B} \supset \bigcup_{i \in \mathbb{N}} \overline{A_{i}}$.
	\end{enumerate}
\end{problem}

\pagebreak

\begin{problem}[8.]
	Is every point of an open set $E \subset \mathbb{R}^{2}$ a limit point of $E$? Answer the same question for closed sets in $\mathbb{R}^{2}$.
\end{problem}

\pagebreak

\begin{problem}[8*.]
	Are there any infinite metric spaces which have no infinite compact subsets?
\end{problem}

\pagebreak

\begin{problem}[9.]
	Let $E^{\circ}$ denote the set of all interior points of a set $E$. $E^{\circ}$ is called the \textit{interior} of $E$.
	\begin{enumerate}
		\item Prove that $E^{\circ}$ is always open.
		\item Prove that $E$ is open if and only if $E^{\circ} = E$.
		\item If $G \subset E$ and $G$ is open, prove that $G \subset E^{\circ}$.
		\item Prove that the complement of $E^{\circ}$ is the closure of the complement of $E$.
		\item Do $E$ and $\overline{E}$ always have the same interiors?
		\item Do $E$ and $E^{\circ}$ always have the same closures?
	\end{enumerate}
\end{problem}

\pagebreak

\begin{problem}[10.]
	Let $X$ be an infinite set. For $p \in X$ and $q \in X$, define
	\[
		d(p, q) = \begin{cases}
			1 & \text{if } p \neq q \\
			0 & \text{if } p = q.
		\end{cases}
	\]
	Prove that this is a metric. Which subsets of the resulting metric spaces are open? Which are closed? Which are compact?
\end{problem}

\pagebreak

\begin{problem}[11.]
	For $x \in \mathbb{R}$ and $y \in \mathbb{R}$, define
	\[
		\begin{split}
			d_{1}(x, y) & = (x - y)^{2} \\
			d_{2}(x, y) & = \sqrt{\lvert x - y \rvert} \\
			d_{3}(x, y) & = \lvert x^{2} - y^{2} \rvert \\
			d_{4}(x, y) & = \lvert x - 2y \rvert \\
			d_{5}(x, y) & = \frac{\lvert x - y \rvert}{1 + \lvert x - y \rvert}.
		\end{split}
	\]
\end{problem}

\pagebreak

\begin{problem}[12.]
	Let $K \subset \mathbb{R}$ consist of $0$ and the numbers $1/n$ for $n \in \mathbb{N}$. Prove that $K$ is compact directly from the definition.
\end{problem}

\pagebreak

\begin{problem}[13.]
	Construct a compact set of real numbers whose limit points form a countable set.
\end{problem}

\pagebreak

\begin{problem}[14.]
	Give an example of an open cover of the segment $(0, 1)$ which has no finite subcover.
\end{problem}

\pagebreak

\begin{problem}[15.]
	Show that Theorem 2.36 (that the intersection of a collection of compact sets with all nonempty finite intersection is nonempty) and its corollary become false (in $\mathbb{R}$ for example) if the word ``compact'' is replaced by ``closed'' or ``bounded''.
\end{problem}

\pagebreak

\begin{problem}[15*.]
	Prove that every open set in $\mathbb{R}$ is the union of an at most countable collection of disjoint segments.
\end{problem}

\pagebreak

\begin{problem}[16.]
	Regard $\mathbb{Q}$, the set of all rational numbers, as a metric space, with $d(p, q) = \lvert p - q \rvert$. Let $E$ be the set of all $p \in \mathbb{Q}$ such that $2 < p^{2} < 3$. Show that $E$ is closed and bounded in $\mathbb{Q}$, but that $E$ is not compact. Is $E$ open in $\mathbb{Q}$?
\end{problem}

\pagebreak

\begin{problem}[17.]
	Let $E$ be the set of all $x \in [0, 1]$ whose decimal expansion contains only the digits $4$ and $7$. Is $E$ countable? Is $E$ dense in $[0, 1]$? Is $E$ compact? Is $E$ perfect?
\end{problem}

\pagebreak

\begin{problem}[18.]
	Is there a nonempty perfect set in $\mathbb{R}$ which contains no rational number?
\end{problem}

\pagebreak

\begin{problem}[19.]
	\begin{enumerate}
		\item If $A$ and $B$ are disjoint closed sets in some metric space $X$, prove that they are separated.
		\item Prove the same for disjoint open sets.
		\item Fix $p \in X$, $\delta > 0$, define $A$ to be the set of all $q \in X$ for which $d(p, q) < \delta$ define $B$ similarly with $>$ in place of $<$. Prove that $A$ and $B$ are separated.
		\item Prove that every connected metric space with at least two points is uncountable.
	\end{enumerate}
\end{problem}

\pagebreak

\begin{problem}[20.]
	Are closures and interiors of connected sets always connected?
\end{problem}

\pagebreak

\begin{problem}[21.]
	Let $A$ and $B$ be separated subsets of some  $\mathbb{R}^{k}$, suppose $\mathbf{a} \in A$, $\mathbf{b} \in B$, and define
	\[
		p(t) = (1 - t)\mathbf{a} + t \mathbf{b}
	\]
	for $t \in \mathbb{R}$. Put $A_{0} = p^{\text{pre}}(A), B_{0} = p^{\text{pre}}(B)$.
	\begin{enumerate}
		\item Prove that $A_{0}$ and $B_{0}$ are separated subsets of $\mathbb{R}$.
		\item Prove that there exists $t_{0} \in (0, 1)$ such that $p(t_{0}) \notin A \cup B$.
		\item Prove that every convex subset of $\mathbb{R}^{k}$ is connected.
	\end{enumerate}
\end{problem}

\pagebreak

\begin{problem}[22.]
	A metric space is called \textit{separable} if it contains a countable dense subset. Show that $\mathbb{R}^{k}$ is separable.
\end{problem}

\pagebreak

\begin{problem}[23.]
	A collection $\{ V_{\alpha} \}$ of open sets of $X$ is said to be a \textit{base} for $X$ if the following is true: for every $x \in X$ and every open set $G \subset X$ such that $x \in G$, we have $x \in V_{\alpha} \subset G$ for some $\alpha$. In other words, every open set in $X$ is the union of a subcollection of $\{ V_{\alpha} \}$. \\

	Prove that every separable metric space has a countable base.
\end{problem}

\pagebreak

\begin{problem}[24.]
	Let $X$ be a metric space in which every infinite subset has a limit point. Prove that $X$ is separable.
\end{problem}

\pagebreak

\begin{problem}[25.]
	Prove that every compact metric space $K$ has a countable base, and that $K$ is therefore separable.
\end{problem}

\pagebreak

\begin{problem}[26.]
	Let $X$ be a metric space in which every infinite subset has a limit point. Prove that $X$ is compact.
\end{problem}

\pagebreak

\begin{problem}[27.]
	Define a point $p$ in a metric space $X$ to be a \textit{condensation point} of a set $E \subset X$ if every neighbourhood of $p$ contains uncountably many points of $E$.

	Suppose $E \subset \mathbb{R}^{k}$, $E$ is uncountable, and let $P$ be the set of all condensation points of $E$. Prove that $P$ is perfect and that at most countably many points of $E$ are not in $P$. In other words, show that $P^{c} \cap E$ is at most countable.
\end{problem}

\pagebreak

\begin{problem}[28.]
	Prove that every closed set in a separable metric space is the union of a (possibly empty) perfect set and a set which is at most countable.
\end{problem}

\pagebreak

\begin{problem}[29.]
	Prove that every open set in $\mathbb{R}$ is the union of an at most countable collection of disjoint segments.
\end{problem}

\pagebreak

\begin{problem}[30.]
	Imitate the proof of Theorem 2.43 to obtain the following result: 

	If $\mathbb{R}^{k} = \bigcup_{n \in \mathbb{N}} F_{n}$ where each $F_{n}$ is a closed subset of  $\mathbb{R}^{k}$, then at least one $F_{n}$ has a nonempty interior. \\
\end{problem}

\end{document}
